\enquestion{7.34}{
    Let $X$ be a geometric random variable with parameter $p$.
    Find the maximum likelihood estimator of $p$ based on a random sample of size $n$.
}
\enanswer{
    $\hat{p} = \frac{1}{\bar{X}}$
}

\enquestion{7.35}{
    Consider the Poisson distribution with parameter $\lambda$.
    Find the maximum likelihood estimator of $\lambda$, based on a random sample of size $n$.
}
\enanswer{
    $\hat{\lambda} = \bar{X}$
}

\enquestion{7.38}{
    Consider the probability density function
    \[
        f(x) = \frac{1}{\theta^2} x e^{-x / \theta}, \quad 0 \le x < \infty, \quad 0 < \theta < \infty
    \]
    Find the maximum likelihood estimator for $\theta$.
}
\enanswer{
    $\hat{\theta} = \frac{\bar{X}}{2}$
}

\enquestion{7.40}{
    Consider the probability density function
    \[
        f(x) = c(1 + \theta x), \quad -1 \le x \le 1
    \]
    \begin{enumerate}[label=(\alph*)]
        \item Find the value of the constant $c$.
        \item Show that $\hat{\theta} = 3\bar{X}$ is an unbiased estimator for $\theta$.
    \end{enumerate}
}
\enanswer{
    \begin{enumerate}[label=(\alph*)]
        \item $c = \frac{1}{2}$
        \item $E(\bar{X}) = E(X) = \frac{1}{3}\theta$. So $E(3\bar{X}) = \theta$.
    \end{enumerate}
}

\enquestion{7.41}{
    The Rayleigh distribution has probability density function
    \[
        f(x) = \frac{x}{\theta} e^{-x^{2} / 2\theta}, \quad x > 0, \quad 0 < \theta < \infty
    \]
    \begin{enumerate}[label=(\alph*)]
        \item It can be shown that $E(X^{2}) = 2\theta$.
        Use this information to construct an unbiased estimator for $\theta$.
        \item Find the maximum likelihood estimator of $\theta$.
        Compare your answer to part (a).
    \end{enumerate}
}
\enanswer{
    \begin{enumerate}[label=(\alph*)]
        \item $\hat{\theta} = \frac{\sum_{i=1}^{n} X_{i}^{2}}{2n}$
        \item $\hat{\theta} = \frac{\sum_{i=1}^{n} X_{i}^{2}}{2n}$, the maximum likelihood estimator is identical to the unbiased estimator in part (a).
    \end{enumerate}
}

\enquestion{7.42}{
    Let $X_{1}, X_{2}, \dots, X_{n}$ be uniformly distributed on the interval $0$ to $a$.
    Recall that the maximum likelihood estimator of $a$ is $\hat{a} = \max(X_{i})$.
    \begin{enumerate}[label=(\alph*)]
        \item Argue intuitively why $\hat{a}$ cannot be an unbiased estimator for $a$.
        \item Suppose that $E(\hat{a}) = na / (n + 1)$.
        Is it reasonable that $\hat{a}$ consistently underestimates $a$?
        Show that the bias in the estimator approaches zero as $n$ gets large.
        \item Propose an unbiased estimator for $a$.
        \item Let $Y=\max(X_i)$.
        Use the fact that $Y \leq y$ if and only if each $X_i \leq y$ to derive the cumulative distribution function of $Y$.
        Then show that the probability density function of $Y$ is
        \[
            f(y) = \begin{cases}
                \frac{ny^{n-1}}{a^n},   & 0\leq y \leq a \\
                0,                      & \text{otherwise}                      
            \end{cases}
        \]  
        Use this result to show that the maximum likelihood estimator for $a$ is biased.
        \item We have two unbiased estimator for $a$: the moment estimator $\hat{a}_1 = 2\bar{X}$ and $\hat{a}_2=[(n+1)/n]\max(X_i)$, where $\max(X_i)$ is the largest observation in a random sample of size $n$.
        It can be shown that $V(\hat{a}_1)=a^2/(3n)$ and that $V(\hat{a}_2)=a^2/[n(n+2)]$.
        Show that if $n>1$, $\hat{a}_2$ is a better estimator than $\hat{a}_1$.
        In what sense is it a better estimator of $a$?
    \end{enumerate}
}
\enanswer{
    \begin{enumerate}[label=(\alph*)]
        \item The largest observation is almost always less than $a$.
        \item Since $\frac{n}{n+1} < 1$, it consistently underestimates $a$. The bias is $\frac{na}{n+1} - a = -\frac{a}{n+1} \to 0$ as $n \to \infty$.
        \item $\hat{a} = \frac{n+1}{n} X_{\max}$
        \item $F(y) = \left(\frac{y}{a}\right)^n \Rightarrow f(y) = \frac{ny^{n-1}}{a^n} \Rightarrow E(Y) = \frac{na}{n+1} \Leftrightarrow E\left(\frac{n+1}{n}Y\right) = a$, showing that $Y = \hat{a}_{\text{MLE}}$ is biased.
        \item Since $n(n+2) > 3n \Leftrightarrow n < 0$ or $n > 1 \Leftrightarrow n > 1$, $V(\hat{a}_2) < V(\hat{a}_1)$ for $n > 1$. Thus, $\hat{a}_2$ is a better (i.e., more precise) estimator than $\hat{a}_1$.
    \end{enumerate}
}